% ---------------------------------------------------------------------------
% preamble.tex — typography for the paper PDFs.
%
% Chosen for a text-heavy scientific paper with Cyrillic in the Russian edition:
%   * main font  P052 (URW Palladio / Palatino clone) — a book face with real
%     Cyrillic coverage, which DejaVu Serif (the previous build) has but which
%     looks like a screen font in print;
%   * the nine glyphs P052 does not carry are mapped explicitly (see below) —
%     measured with fontTools, not guessed: the missing set is
%     U+207B ⁻, U+2080 ₀, U+2081 ₁, U+2082 ₂, U+208A ₊, U+209B ₛ, U+0307 ̇,
%     U+27F9 ⟹, U+2293 ⊓.
% ---------------------------------------------------------------------------

\usepackage{fontspec}
\usepackage{unicode-math}

% --- fonts -----------------------------------------------------------------
\setmainfont{P052}[
  Extension      = .otf,
  UprightFont    = *-Roman,
  BoldFont       = *-Bold,
  ItalicFont     = *-Italic,
  BoldItalicFont = *-BoldItalic,
  Numbers        = OldStyle
]
\setsansfont{TeX Gyre Heros}
\setmonofont{DejaVu Sans Mono}[Scale=0.86]
\setmathfont{Latin Modern Math}

% --- the glyphs the main font lacks ---------------------------------------
% Measured with fontTools, not guessed: P052 is missing exactly
%   U+207B ⁻, U+2080 ₀, U+2081 ₁, U+2082 ₂, U+208A ₊, U+209B ₛ, U+27F9 ⟹, U+2293 ⊓
% (and U+0307, handled in the build script).
%
% These are mapped to a FALLBACK FONT THAT HAS THEM (DejaVu Serif) rather than to
% a LaTeX construction like \textsuperscript{\textminus}. The difference matters:
% with the construction the PDF renders correctly but the TEXT LAYER loses the
% character — copy-paste, search and text extraction all degrade, and a reader
% searching for "10⁻²³" would not find it. With the fallback font the actual
% Unicode character survives into the PDF, which is what verify_pdfs.py checks.
\usepackage{newunicodechar}
\newfontfamily\glyphfallback{DejaVu Serif}
\newunicodechar{⁻}{{\glyphfallback ⁻}}
\newunicodechar{⁺}{{\glyphfallback ⁺}}
\newunicodechar{⁰}{{\glyphfallback ⁰}}
\newunicodechar{¹}{{\glyphfallback ¹}}
\newunicodechar{²}{{\glyphfallback ²}}
\newunicodechar{³}{{\glyphfallback ³}}
\newunicodechar{⁴}{{\glyphfallback ⁴}}
\newunicodechar{⁵}{{\glyphfallback ⁵}}
\newunicodechar{⁶}{{\glyphfallback ⁶}}
\newunicodechar{⁷}{{\glyphfallback ⁷}}
\newunicodechar{⁸}{{\glyphfallback ⁸}}
\newunicodechar{⁹}{{\glyphfallback ⁹}}
\newunicodechar{₀}{{\glyphfallback ₀}}
\newunicodechar{₁}{{\glyphfallback ₁}}
\newunicodechar{₂}{{\glyphfallback ₂}}
\newunicodechar{₃}{{\glyphfallback ₃}}
\newunicodechar{₄}{{\glyphfallback ₄}}
\newunicodechar{₊}{{\glyphfallback ₊}}
\newunicodechar{ₛ}{{\glyphfallback ₛ}}
\newunicodechar{⟹}{{\glyphfallback ⟹}}
\newunicodechar{⊓}{{\glyphfallback ⊓}}
\newunicodechar{⊔}{{\glyphfallback ⊔}}

% The combining dot above (U+0307) cannot be redefined by \newunicodechar on its
% own — it is a combining mark and must stay in the SAME font as the base letter
% to compose. So the build script maps the pair "tau + U+0307" to \tauDot, which
% sets BOTH characters in the fallback font: they compose correctly, and the text
% layer keeps U+03C4 U+0307 (checked by verify_pdfs.py).
\newcommand{\tauDot}{{\glyphfallback τ̇}}

% --- page and text ---------------------------------------------------------
\usepackage[a4paper, margin=2.4cm, top=2.6cm, bottom=2.4cm]{geometry}
% NOTE: pandoc's default template already loads microtype, so it is NOT loaded
% here — a second \usepackage with different options is an "Option clash" error.
% Protrusion is disabled through pandoc's own -V microtypeoptions instead (see
% build_paper.sh): microtype's character protrusion lets a glyph extend ~2pt into
% the margin, which is standard typography and invisible — but it makes "is every
% word inside the text block?" unanswerable, and the owner asked exactly that
% question. Expansion and kerning are kept; only protrusion is off, so the frame
% check in verify_pdfs.py is exact rather than tolerance-based.
\usepackage{setspace}
\setstretch{1.06}
\usepackage{parskip}
\setlength{\parskip}{0.55em}
\usepackage{ragged2e}

% --- nothing may leave the text block --------------------------------------
% The owner found a formula running past the right margin. TeX's default
% tolerance prefers an overfull box to a loose line; that is the wrong trade for
% a paper, so the tolerance is raised and an emergency stretch is allowed.
\usepackage{seqsplit}   % lets a long unbreakable token be split with \seqsplit
\usepackage{xurl}       % lets a long URL break at any character
\Urlmuskip=0mu plus 1mu
\setlength{\emergencystretch}{3em}
\tolerance=3000
\hbadness=10000
\vbadness=10000
\sloppy
% long technical words that must be allowed to break
\hyphenation{Schwarz-schild-me-tric Schwarz-schild iden-ti-fi-a-bil-ity
             iden-ti-fi-able IMR-Phe-nom-T Math-lib}

% --- headings --------------------------------------------------------------
\usepackage{titlesec}
\usepackage{xcolor}
\definecolor{inkblue}{HTML}{1B3A6B}
\definecolor{inkred}{HTML}{A8322D}
\definecolor{inkgold}{HTML}{8A6D1F}
\definecolor{rulegrey}{HTML}{B9BFC9}

\titleformat{\section}
  {\normalfont\Large\bfseries\color{inkblue}}{\thesection}{0.7em}{}
\titleformat{\subsection}
  {\normalfont\large\bfseries\color{inkblue!85!black}}{\thesubsection}{0.6em}{}
\titleformat{\subsubsection}
  {\normalfont\normalsize\bfseries\color{inkgold}}{\thesubsubsection}{0.5em}{}
\titlespacing*{\section}{0pt}{1.7em}{0.7em}
\titlespacing*{\subsection}{0pt}{1.3em}{0.5em}
\titlespacing*{\subsubsection}{0pt}{1.1em}{0.4em}

% a hairline under section titles
\newcommand{\sectionrule}{\vspace{-0.35em}{\color{rulegrey}\hrule height 0.4pt}\vspace{0.35em}}
\titleformat{\section}
  {\normalfont\Large\bfseries\color{inkblue}}{\thesection}{0.7em}{}[\sectionrule]

% --- tables ----------------------------------------------------------------
\usepackage{booktabs}
\usepackage{array}
\usepackage{longtable}
\renewcommand{\arraystretch}{1.18}
\setlength{\tabcolsep}{5pt}

% --- code ------------------------------------------------------------------
\usepackage{fancyvrb}
\definecolor{codebg}{HTML}{F5F6F8}
\definecolor{codeink}{HTML}{243447}

% --- links -----------------------------------------------------------------
\usepackage{hyperref}
\hypersetup{
  colorlinks = true,
  linkcolor  = inkblue,
  citecolor  = inkblue,
  urlcolor   = inkblue,
  filecolor  = inkblue,
  pdftitle   = {Is the Müller–Maguire "now"-theory prediction falsifiable on GW150914?},
  pdfauthor  = {Aiodam; Oxunjon Ubaydullayev},
  pdfsubject = {Formal identifiability audit with a numerical check on real LIGO strain}
}

% --- headers and footers ---------------------------------------------------
\usepackage{fancyhdr}
\pagestyle{fancy}
\fancyhf{}
\renewcommand{\headrulewidth}{0.3pt}
\renewcommand{\footrulewidth}{0pt}
\fancyhead[L]{\footnotesize\color{inkblue}\itshape Identifiability of the “now”-theory delay on GW150914}
\fancyhead[R]{\footnotesize\color{inkblue}\thepage}
\fancyfoot[C]{\footnotesize\color{inkgold} Aiodam \& O. Ubaydullayev — 2026}

% --- title block -----------------------------------------------------------
% NOTE: \preauthor / \postauthor are deliberately NOT redefined. titling's
% default wraps the author list in a tabular (which is how \and is rendered);
% replacing it with a plain \begin{center} breaks \maketitle with
% "! Misplaced \crcr." — measured, not guessed.
\usepackage{titling}
\pretitle{\begin{center}\vspace{-1.2em}\color{inkblue}\LARGE\bfseries}
\posttitle{\par\end{center}\vspace{0.2em}}
\predate{\begin{center}\small\color{inkgold}}
\postdate{\par\end{center}\vspace{0.6em}}

% --- abstract --------------------------------------------------------------
\renewenvironment{abstract}{%
  \begin{center}%
    \begin{minipage}{0.92\textwidth}%
      \setstretch{1.0}\small%
      \noindent{\color{inkblue}\bfseries Abstract}\par\vspace{0.35em}%
      \hrule height 0.4pt \vspace{0.5em}%
}{%
      \vspace{0.5em}\hrule height 0.4pt%
    \end{minipage}%
  \end{center}%
  \vspace{0.8em}%
}

% --- figures ---------------------------------------------------------------
\usepackage{graphicx}
% Figures are placed EXACTLY where the text puts them ([H]). That is not a
% stylistic choice: with the default deferred placement LaTeX may pile two
% floats onto one page, and when it cannot fit them it emits an "Overfull \vbox
% ... has occurred while \output is active" and the second caption is set past
% the bottom of the page — ink outside the text block, which is exactly what
% verify_frames.py exists to catch. Measured: with deferred placement the EN
% paper overflows on two pages (188 px past the block, to the physical edge);
% with [H] the overflow is zero and the ink stays inside. The trade is that a
% figure can leave a gap at the bottom of a page; correctness first.
\usepackage{float}
\floatplacement{figure}{H}
\usepackage{caption}
\captionsetup{font=small, labelfont={bf,color=inkblue}, labelsep=period,
              justification=raggedright, singlelinecheck=false}
\setlength{\abovecaptionskip}{0.55em}
\setlength{\belowcaptionskip}{0.25em}

% --- lists -----------------------------------------------------------------
\usepackage{enumitem}
\setlist[itemize]{topsep=0.35em, itemsep=0.15em, parsep=0.1em, leftmargin=1.3em}
\setlist[enumerate]{topsep=0.35em, itemsep=0.15em, parsep=0.1em, leftmargin=1.6em}

% --- block quotes (the paper uses them for the formal statements) ----------
\usepackage{mdframed}
\renewenvironment{quote}{%
  \begin{mdframed}[linewidth=0.6pt, linecolor=rulegrey, backgroundcolor=codebg!60,
                   innertopmargin=0.6em, innerbottommargin=0.6em,
                   innerleftmargin=0.9em, innerrightmargin=0.9em,
                   skipabove=0.6em, skipbelow=0.6em]%
  \small%
}{%
  \end{mdframed}%
}

% --- widow/orphan control --------------------------------------------------
\widowpenalty=10000
\clubpenalty=10000